	% ==========================================
% BAB VI EVALUASI
% ==========================================
\cleardoublepage
\chapter{EVALUASI}
\label{chap:evaluasi}
asdfasdf
\section{Metode Evaluasi}
Berdasarkan 
\section{Hasil Evaluasi}
\subsection{Model}
\subsubsection{CLM}
\subsection{Analisis Kritis Hasil Evaluasi Model}

Evaluasi terhadap sembilan skenario eksperimen menunjukkan bahwa metrik akurasi dan \textit{composite score} yang tinggi tidak selalu merepresentasikan kemampuan model yang sesungguhnya, sehingga diperlukan pembahasan lebih mendalam terhadap beberapa temuan berikut.

\begin{enumerate}
	
	\item \textbf{Kolaps Kelas (\textit{Class Collapse}) dan Pengecualian pada Skema WN--\textit{all}}\\
	Enam dari sembilan eksperimen menghasilkan model terbaik dengan sensitivitas 100\% namun spesifisitas 0\%, yang berarti model selalu memprediksi kelas tumor tanpa benar-benar mempelajari pola pembeda antar kelas. Nilai akurasi dan \textit{composite score} yang tampak tinggi pada kondisi ini bersifat menyesatkan karena hanya merefleksikan proporsi kelas mayoritas pada data uji, dan mengindikasikan ketidakseimbangan kelas (\textit{class imbalance}) yang belum ditangani secara memadai. Satu-satunya pengecualian adalah kombinasi \textit{wavelet denoising} dan normalisasi intensitas (WN) pada skema \textit{all}, yang menghasilkan sensitivitas dan spesifisitas relatif seimbang (62,03\% dan 38,10\%). Meskipun akurasi mentahnya lebih rendah, hasil ini lebih dapat dipercaya secara klinis karena model benar-benar menangkap sebagian sinyal pembeda antar kelas.
	
	\item \textbf{Instabilitas Pelatihan dan \textit{Domain Shift} pada Skema Penggabungan Dataset}\\
	Skema penggabungan dataset kedua secara penuh ke data latih, validasi, dan uji (\textit{all}) menunjukkan pelatihan yang paling tidak stabil, dengan \textit{cross-entropy loss} yang tertahan di sekitar nilai \textit{random-guess} ($\approx \ln 2 = 0{,}693$) atau berosilasi tanpa tren menurun. Sebaliknya, skema \textit{test} dan \textit{val\_test} menunjukkan penurunan \textit{loss} yang lebih mulus dengan akurasi latih tinggi (90--96\%), namun beberapa konfigurasi mencatat selisih besar antara akurasi validasi dan uji (mis. skema \textit{gc\_test}, \textit{dropout} 0,05: validasi 70,16\% vs uji 29,41\%). Kedua temuan ini mengindikasikan bahwa dataset kedua memiliki distribusi fitur yang cukup berbeda dari dataset pertama, sehingga baik pencampuran penuh maupun kemunculannya hanya pada data uji dapat mengganggu proses pembelajaran atau generalisasi model.
	
	\item \textbf{Pengaruh \textit{Learning Rate} Tinggi dan Iterasi Berlebih}\\
	Konfigurasi dengan \textit{learning rate} 0,001 serta total iterasi 10.000 secara konsisten menghasilkan \textit{composite score} yang lebih rendah dibandingkan konfigurasi \textit{baseline} (\textit{learning rate} 0,00025, 5.000 iterasi) pada hampir seluruh skema, menunjukkan model cukup sensitif terhadap \textit{learning rate} yang terlalu besar dan berisiko tidak stabil atau \textit{overfitting} bila iterasi ditambah tanpa regularisasi tambahan.
	
	\item \textbf{Kesenjangan terhadap Studi Acuan CLM \parencite{wozniak2023}}\\
	Dibandingkan dengan Wo\'{z}niak dkk. (2023), yang melaporkan akurasi 95--100\% dengan sensitivitas dan spesifisitas seimbang di atas 88\%, hasil pada penelitian ini masih jauh di bawahnya bahkan pada skenario terbaik (WN--\textit{all}: sensitivitas 62,03\%, spesifisitas 38,10\%). Kesenjangan ini kemungkinan besar disebabkan oleh perbedaan skala anggaran komputasi---studi acuan melatih model hingga 120.000--12.000.000 iterasi dengan implementasi \textit{custom} \textit{multi-threaded} (mekanisme \textit{palette archiving} dan \textit{reshuffling}), sedangkan penelitian ini hanya menggunakan 3.000--10.000 iterasi dengan Adam standar tanpa mekanisme arsip-palet yang setara. Selain itu, dataset kedua pada studi acuan digunakan untuk klasifikasi tiga kelas, berbeda dengan skema biner pada penelitian ini, sehingga perbandingan akurasi antara kedua penelitian tidak sepenuhnya \textit{apple-to-apple}.
	
\end{enumerate}

\subsubsection{Explainable CNN}
\subsection{Masalah Model yang Hanya Menebak Satu Kelas}
Pada sebagian kombinasi \textit{hyperparameter}, model tidak benar-benar mempelajari perbedaan kelas tumor dan bukan tumor, melainkan hanya menebak satu kelas untuk seluruh data uji. Hal ini terlihat dari nilai \textit{specificity} yang tepat 0\% (model selalu menebak kelas tumor) atau 100\% (model selalu menebak kelas bukan tumor), sehingga akurasi yang terlihat tinggi pada kasus ini tidak mencerminkan kemampuan model yang sesungguhnya. Dari 27 kombinasi \textit{hyperparameter} pada tiap skenario, persentase model yang mengalami masalah ini berkisar antara 25,9\% (WN dengan \textit{val\_test}) hingga 85,2\% (tanpa \textit{preprocessing} dan GC dengan \textit{val\_test}). Oleh karena itu, nilai \textit{specificity} dan \textit{F1-Score} perlu selalu diperhatikan bersama akurasi dalam menilai model.

\subsection{Indikasi \textit{Overfitting} pada \textit{Learning Curve}}
Pada hampir seluruh grafik \textit{learning curve} model terbaik di tiap skenario, akurasi data \textit{training} naik cepat hingga mendekati 100\% dalam kurang dari 10 \textit{epoch}, sementara akurasi data validasi tidak ikut naik dan cenderung tidak stabil, bahkan \textit{loss} validasi meningkat setelah beberapa \textit{epoch} awal. Pola ini menunjukkan \textit{overfitting}, yaitu kondisi model menghafal data \textit{training} alih-alih mempelajari pola umum, yang kemungkinan disebabkan oleh jumlah data \textit{training} yang terbatas dan nilai \textit{dropout} (0,1--0,3) yang belum cukup untuk mencegahnya.

\subsection{Pengaruh \textit{Preprocessing} terhadap Kemampuan Generalisasi}
Pada skenario \textit{test}, yaitu model dilatih dari \textit{Dataset} 1 lalu diuji langsung pada \textit{Dataset} 2 yang belum pernah dilihat sebelumnya, tanpa \textit{preprocessing} tidak ada satupun dari 27 kombinasi \textit{hyperparameter} yang menghasilkan \textit{specificity} di atas 0\%. Setelah \textit{preprocessing} GC atau WN diterapkan, ditemukan kombinasi yang menghasilkan \textit{specificity} jauh lebih baik (GC: 81\%, WN: 50\%), menunjukkan bahwa \textit{denoising} dan normalisasi intensitas berperan penting dalam menjembatani perbedaan karakteristik antar \textit{dataset}.

\subsection{Pengaruh Strategi Penggabungan \textit{Dataset}}
Tidak ditemukan satu strategi penggabungan \textit{dataset} (\textit{all}, \textit{test}, atau \textit{val\_test}) yang selalu terbaik pada semua kondisi. Tanpa \textit{preprocessing} dan dengan WN, strategi \textit{val\_test} memberikan hasil terbaik, sedangkan dengan GC strategi \textit{all} justru lebih unggul. Hal ini menunjukkan bahwa efektivitas strategi penggabungan \textit{dataset} tampak bergantung pada jenis \textit{preprocessing} yang digunakan, sehingga belum dapat ditarik kesimpulan tunggal yang berlaku umum.

\subsection{Pengaruh \textit{Hyperparameter}}
\textit{Learning rate} 0,0001 dan 0,001 cenderung memberikan performa lebih baik dan stabil dibanding 0,01, terutama pada skenario dengan \textit{preprocessing} WN. \textit{Dropout} 0,1 cenderung unggul tipis pada skenario dengan performa terbaik, sedangkan \textit{batch size} tidak menunjukkan pola pengaruh yang konsisten pada semua skenario. Karena setiap kombinasi hanya dijalankan satu kali tanpa pengulangan, kesimpulan mengenai pengaruh \textit{hyperparameter} ini bersifat kecenderungan umum dan belum sepenuhnya konklusif.
\subsubsection{EfficientNetB0}
\begin{enumerate}
	\item \textbf{Overfitting yang Konsisten dan Berat}:
	Pada seluruh skema eksperimen, akurasi \textit{training} stabil mencapai 95--99\%, sementara akurasi \textit{validation} berfluktuasi jauh lebih rendah (55--85\%) dengan lonjakan \textit{validation loss} yang tajam, terutama setelah fase \textit{fine-tuning} dimulai. Pola ini konsisten pada hampir semua kombinasi \textit{hyperparameter} dan varian \textit{preprocessing}, sehingga \textit{overfitting} yang terjadi bersifat sistemik pada pipeline, bukan akibat pemilihan \textit{hyperparameter} tertentu, dan mengindikasikan kapasitas model relatif besar dibandingkan ukuran dan variasi data yang tersedia.
	
	\item \textbf{Efek Strategi Penggabungan \textit{Dataset}}:
	Ketiga strategi penggabungan dataset kedua (\textit{all}, \textit{test}, \textit{val\_test}) tidak menunjukkan satu strategi yang konsisten superior, dengan selisih rata-rata akurasi yang relatif kecil (86,5\%--87,8\%). Strategi \textit{all} membuat distribusi data uji lebih mirip data latih namun mengurangi independensi data uji, sedangkan \textit{test} dan \textit{val\_test} menjadikan dataset kedua benar-benar \textit{unseen} sehingga lebih mencerminkan generalisasi lintas-domain, tetapi dengan variansi metrik yang jauh lebih tinggi (std akurasi >7\% vs sekitar 3\% pada \textit{all}). Pemilihan strategi ini merupakan \textit{trade-off} antara validitas evaluasi dan stabilitas metrik.
	
	\item \textbf{\textit{Trade-off} Sensitivitas vs Spesifisitas}:
	Pada hampir seluruh model terbaik, nilai \textit{recall} secara konsisten lebih tinggi dibandingkan \textit{specificity} (misalnya 95,88\% berbanding 76,19\%), menunjukkan model lebih baik mendeteksi kasus positif dibandingkan mengenali kasus negatif, sehingga berpotensi menghasilkan lebih banyak \textit{false positive}. Bias ke arah sensitivitas ini relatif dapat diterima dalam deteksi tumor karena \textit{false negative} dianggap lebih berbahaya secara klinis, namun spesifisitas yang berada di kisaran 76--80\% pada beberapa skema tetap merupakan keterbatasan yang berimplikasi pada meningkatnya kasus yang memerlukan pemeriksaan lanjutan.
\end{enumerate}
\subsection{XAI}
\subsubsection{Kuantitatif}
\subsubsection{Kualitatif}
\section{Pembahasan Hasil Evaluasi}